% Complete proof; independently audited by root, B, and C, with primary sources checked.
\subsection{A quadratic presentation of the tame boundary flags}

\begin{lemma}[Proved: linear resolutions of an entire bounded-slope flag]
\label{D:filtered-linear-curve-resolution}
Fix a smooth projective curve $C\subset\mathbf P^N_k$ and real numbers
$a\le b$.  There are integers $m,K$ with the following property.
For every flag of vector subbundles
\[
        0=V_0\subset V_1\subset\cdots\subset V_l=V,
            \qquad l\le r=\operatorname{rk}V,
\]
whose nonzero grades $G_j=V_j/V_{j-1}$ have all slopes in $[a,b]$,
the flag is the sheafification of a filtered graded free complex
over $P=k[x_0,\ldots,x_N]$.  This complex has length at most $N+1$,
its term in homological degree $i$ is a sum of at most $Kr$
copies of $P(-m-i)$, and its differential is a matrix of linear
forms preserving coordinate flags in all terms.  The associated
graded complex resolves the corresponding truncated section
modules of the $G_j$.

For fixed ranks and multiplicities, these filtered complexes are
parameterized by a closed affine scheme with $O(r^2)$ variables
and $O(r^2)$ quadratic equations.  Exactness, flatness and the
subbundle conditions are imposed by passing to an open subset
after identifying the total cohomology with a prescribed bundle
on $C$.  The number of multiplicity arrays needed is at most
$\exp(Kr)$.
\end{lemma}

\begin{proof}
Choose $m$ so large that
$a+(m-1)\deg\mathcal O_C(1)>2g(C)-2$.
Every grade, every partial subbundle and every partial quotient
then has vanishing $H^1$ in all twists of degree at least $m-1$.
This follows from the slope bounds and their preservation under
extensions.  A fixed further increase of $m$ ensures global
generation, by applying the same vanishing after subtracting
one point.

For such a bundle $W$ put
\[
          M(W)=\bigoplus_{d\ge m}H^0(C,W(d)).
\]
It is a finitely generated graded $P$-module, with no components
below $m$.  The local-cohomology comparison with its associated
sheaf shows that $H^0_{\mathfrak m}M(W)=0$, that
$H^1_{\mathfrak m}M(W)$ vanishes in degrees at least $m$,
and that $H^2_{\mathfrak m}M(W)$ vanishes in degrees at
least $m-1$.  Higher local cohomology vanishes because the
sheaf is supported on a curve.  Hence $M(W)$ is $m$-regular.
The equivalence of regularity with linear truncation
\citep[Proposition~1.1 and Theorem~1.2, pp.~95--97,
including the proof]{EisenbudGoto1984D}
gives an $m$-linear graded free resolution.
Hilbert's syzygy theorem bounds its length by $N+1$.

Its Betti numbers are linear in the rank bound.  Indeed, for
$d\ge m$ Riemann--Roch gives
\[
       h^0(C,W(d))
        =\operatorname{rk}(W)\deg\mathcal O_C(1)\,d
                      +\deg W+\operatorname{rk}(W)(1-g).
\]
The coefficients of
$(1-t)^{N+1}\sum_{d\ge m}h^0(C,W(d))t^d$
are therefore bounded by $K\operatorname{rk}(W)$.
For an $m$-linear resolution, each homological degree
occurs in its own internal degree $m+i$, so these
coefficients, with alternating signs, are precisely
the Betti numbers.

The $H^1$ vanishing makes
\[
 0\to M(V_{j-1})\to M(V_j)\to M(G_j)\to0
\]
exact.  Apply the graded horseshoe construction successively
to the $m$-linear resolutions of the grades.  In degree $i$
all summands have the same shift $-m-i$; all off-diagonal
differentials consequently still have entries of degree one.
The resulting resolution is filtered by coordinate summands,
and its associated graded is the chosen sum of grade resolutions.
Its total multiplicities are bounded by
$K\sum_j\operatorname{rk}G_j=Kr$, rather than by $Kr^2$.

Choose the coordinate flags once the grade multiplicities
are fixed.  Preservation of them says that specified blocks
of each differential are zero.  The remaining blocks have
$O(r^2)$ scalar coefficients.  The equations $d^2=0$ have
$O(r^2)$ scalar coefficients and degree two, since the
complex length and the number of variables of $P$ are fixed.
The loci where the total and associated-graded complexes
have only base-flat degree-zero cohomology, and where
those cohomology sheaves are vector bundles of the
prescribed ranks on $C$, are open by the relative
perfect-complex truncation criterion already used in
Lemma~\ref{D:polynomial-monad-diagrams}.
The exact filtered complex then gives actual subbundles
and locally free quotients over every Artinian parameter
ring, not only at its reduced points.

For each of the fixed $N+2$ degrees the grade multiplicities
are nonnegative integers with sum at most $Kr$, in at
most $r$ ordered positions.  The number of such arrays
is bounded by a fixed power of
$\binom{Kr+r}{r}$, hence by $\exp(K'r)$.
\end{proof}

\begin{lemma}[Proved: quadratic witnesses for filtered diagrams]
\label{D:filtered-polynomial-diagrams}
The conclusion of Lemma~\ref{D:polynomial-monad-diagrams}
remains valid for split-filtered complexes of the preceding
lemma, provided the total multiplicities in each degree are
$O(r)$ and the complex lengths and bundle types are fixed.
The number of filtration steps may be at most $r$.
Morphisms required to preserve the filtration, and equalities
of these morphisms, have witnesses in $O(r^2)$ variables
and equations of bounded degree.
\end{lemma}

\begin{proof}
We use the same finite affine cover and the same fixed
contractions of its Cech Hom complexes as in
Lemma~\ref{D:polynomial-monad-diagrams}.  In every
matrix space retain only the blocks allowed by the
two coordinate flags.  These form the filtered Hom
complex.  The fixed Cech differential and its
chosen contraction act on the fixed bundle factor
and therefore preserve this block condition.
The variable complex differentials preserve the
flags as well.

The perturbation still raises the complex-Hom
column.  Its nilpotence length is bounded by the
fixed lengths of the complexes, not by the number
of grades.  Thus the finite geometric-series
elimination produces polynomial formulas of the
same bounded degree.  The dimensions of all
remaining spaces are at most those of the
unfiltered Hom spaces, namely $O(r^2)$.
Cycles, boundaries witnessing equality, and
composition are treated exactly as before.

For clarity, this computes the derived morphisms
of the exact category of filtered sheaves with
strict exact sequences.  Locally, a split-filtered
free term is a sum of the elementary filtered
free modules whose filtration starts at one
specified step.  These are projective for the
strict exact structure.  Consequently the
filtered locally free complexes and the Cech
resolution compute the required morphisms.
An associated-graded quasi-isomorphism is the
weak equivalence used here.  We are not replacing
the category by its Ext algebra or making a
growing-length iterated-extension construction.
All contraction identities are over $k$ and
remain valid after tensoring by any parameter
algebra, including a nonreduced Artinian one.
\end{proof}

\begin{theorem}[Proved: quadratic defining-characteristic growth with original tame boundary]
\label{D:tame-global-quadratic-representation-bound}
Let $S$ be a normal proper connected surface over an
algebraically closed field of characteristic $p>0$,
let $T\subsetneq S$ be closed, and let $Q(S,T)$ require
tameness only at the original codimension-one components
of $T$.  There is a constant $C_{S,T}$ such that
\[
 \#\bigl(
  \operatorname{Hom}_{\mathrm{cont}}
       (Q(S,T),\operatorname{GL}_r(\mathbf F_{p^f}))
       /\operatorname{GL}_r(\mathbf F_{p^f})\bigr)
          \le p^{C_{S,T}fr^2}                           \tag{D.TQ.1}
\]
for every $r,f\ge1$.
This is a representation-count bound, not a claim
that the full group is finitely generated.
\end{theorem}

\begin{proof}
Enlarge the finite omitted set as in
Theorem~\ref{D:tame-parabolic-fixed-rank-growth};
the desired group is a quotient of that larger-open
group, so it suffices to count for the latter.
Fix the smooth projective resolution $Y$, the positive
exceptional cycle $Z$, and the smooth disjoint strict
boundary curves $C_i$ of
Lemma~\ref{D:saturated-boundary-two-sided-slopes}.
Put $Y^\circ=Y-\operatorname{Supp}Z$ and
$C=\sum_iC_i$.  All embeddings and ample twists
below are fixed.

\emph{Linear sizes of all the objects.}
For the coefficient-Frobenius stages choose one
minimal extension $F_a$ on $Y$.  The canonical
Euler estimate and the relative Euler estimate
give $\chi(F_a)=O(r)$; their first Chern classes
have intersection $O(r)$ with fixed divisors.
The forward and inverse pole comparisons give
uniform upper slopes for $F_a$ and its dual,
both on $Y$ and on a fixed smooth ample curve.
Consequently
\[
                    h^j(Y,F_a(t))\le C_t r
\]
for every fixed twist and $0\le j\le2$, by the
same restriction, duality and Riemann--Roch proof
as in Theorem~\ref{D:global-quadratic-representation-bound}.
This also applies after one Frobenius pullback
and after tensoring by any of finitely many fixed
bundles.  Only one $p$-Frobenius step is used.

The full saturated flag on each $F_a|_{C_i}$ has
all graded slopes in one fixed two-sided interval
by Lemma~\ref{D:saturated-boundary-two-sided-slopes}.
Lemma~\ref{D:filtered-linear-curve-resolution}
therefore encodes its entire flag by one filtered
linear complex with total term sizes $O(r)$.
It does not introduce $r$ independent rank-$r$
surface lattices.

\emph{A fixed number of digit transforms.}
Fix a one-step digit type from the preceding
fixed-rank theorem.  Its digits are nondecreasing,
belong to $\{0,\ldots,p-1\}$, and its permutation
preserves source order inside each equal-digit block.
At one boundary curve we write the saturated flag in its
decreasing convention, so that the suffix of grades with
digit at least $h$ is a vector subbundle, denoted $V_{\ge h}$.
This reverses the indexing convention of the increasing flag
in Lemma~\ref{D:filtered-linear-curve-resolution}; it changes
no object or slope bound.  There are only $p-1$ relevant thresholds.
Define the elementary transform
\[
 E_h=\ker\bigl(F_a\longrightarrow
             i_*(F_a|_{C_i}/V_{\ge h})\bigr).
                                                        \tag{D.TQ.2}
\]
It is a vector bundle.  In an adapted local basis,
its columns are the original columns for digit
at least $h$, and $t$ times the remaining columns,
where $t$ is an equation of $C_i$.
Put $H=F_Y^*F_a$ and
\[
 L_c=H+\sum_{h=1}^{p-1}F_Y^*E_h(hC_i)
                      \ \subset\ H((p-1)C_i).           \tag{D.TQ.3}
\]
Its local basis is $t^{-c_j}F^*v_j$, so it too is
a vector bundle.  For the fixed finite collection
of disjoint $C_i$, perform these constructions
independently; the total number of auxiliary
bundles per stage is bounded solely by $p$ and
the fixed number of boundary components.

All these auxiliary objects have the same linear
cohomology bounds.  In (D.TQ.2) the quotient is
a curve bundle with fixed slope bounds and rank
at most $r$.  Its degree and Euler characteristic
are $O(r)$; the kernel is bounded between
$F_a(-C_i)$ and $F_a$.  Riemann--Roch and these
inclusions bound its cohomology in all fixed
twists.  For its one-step Frobenius pullback,
the first Chern character scales by $p$ and the
second by $p^2$; the latter is already $O(r)$.
For (D.TQ.3), the quotient of $H((p-1)C_i)$
by $L_c$ has a filtration of fixed length at most
$p-1$ whose curve grades are Frobenius pulls
of whole flag grades with fixed normal twists.
Their slopes, degrees and cohomology are bounded
linearly in their ranks.  The same conclusions
follow for $L_c$.  Moreover, the inclusions
$H\subset L_c\subset H((p-1)C)$ restrict generically injectively
to every exceptional component.  They give fixed two-sided
exceptional slope windows for $L_c$, because those for $H$
follow from one-step Frobenius and the divisor $(p-1)C$ is fixed.
Thus the exceptional pole lemma applies to both normalized
one-step arrows between $L_c$ and $F_{a+1}$ with one fixed
constant $B$, independently of rank, denominator, and stage.
Beilinson monads for these finitely many objects thus have fixed types and
lengths, with total multiplicities $O(r)$.

\emph{The canonical transformed flag.}
Let $W_h$ be the quotient consisting of the
consecutive grades with digit exactly $h$,
including their full internal flag.  The local
basis calculation gives a canonical decomposition
\[
          L_c|_{C_i}\simeq
          \bigoplus_{h=0}^{p-1}
                    F_{C_i}^*W_h\otimes N_i^h,
           \qquad N_i=\mathcal O_Y(C_i)|_{C_i}.           \tag{D.TQ.4}
\]
This is not an independently chosen splitting.
For $h>0$, the natural map
$F_Y^*E_h(hC_i)|_{C_i}\to L_c|_{C_i}$
factors through the natural surjection onto
$F_{C_i}^*W_h\otimes N_i^h$ and identifies it
with its summand.  For $h=0$, use $H|_{C_i}$
and its quotient onto $F_{C_i}^*W_0$.
In the local basis, columns of larger digit
vanish after restriction, columns of smaller
digit acquire a positive $t$-power, and exactly
the columns of digit $h$ remain.  These
factorizations characterize (D.TQ.4) uniquely.
They are preserved by every adapted change of
basis and commute with arbitrary parameter
base change.

Within each of its at most $p$ summands retain
the Frobenius flag.  The prescribed permutation
shuffles the summands and their subflags while
preserving order within each summand.  This
gives the intrinsic transformed full flag on
$L_c|_{C_i}$.  It is the flag used in the exact
ceiling-valuations criterion of
Theorem~\ref{D:tame-parabolic-fixed-rank-growth}.

\emph{Bounded-degree witness schemes.}
We now spell out how the descriptions above
produce polynomial witnesses, including at
nonreduced parameter points.
Use the Beilinson monads of
Theorem~\ref{D:global-quadratic-representation-bound}
for the surface bundles and auxiliary sheaves,
and the filtered linear complexes of
Lemma~\ref{D:filtered-linear-curve-resolution}
for the boundary flags and their Frobenius
digit-block flags.  Every total term size is
$O(r)$ and all complex lengths and fixed bundle
types are bounded.

Identify the total boundary complex with
$F_a|_{C_i}$ using the fixed resolution of
$\mathcal O_{C_i}$ and the same degree-zero
truncation construction used in the preceding
quadratic theorem.  Require the total and
graded complexes to have only base-flat
degree-zero cohomology, locally free on $C_i$
of the prescribed ranks.  Because the total
cohomology is the actual restriction, its
subquotients are annihilated by the fixed ideal
of $C_i$ over every parameter algebra.
These are open conditions on the perfect
complex families after the derived support
and identification equations have been imposed.

The kernels in (D.TQ.2) are encoded by the
derived exact-sequence diagram, with their
auxiliary monads and a surjective quotient map.
To encode (D.TQ.3), use the natural map from
$H\oplus\bigoplus_hF_Y^*E_h(hC_i)$ to
$H((p-1)C_i)$, introduce its bundle image $L_c$
and its coherent quotient, and impose the
factorization, surjectivity onto the image,
and the exact quotient sequence.  Cone
quasi-isomorphisms and open flatness/truncation
conditions express exactly this diagram.
There are only a fixed number of nodes.
The actual local diagonal description above
shows that every genuine representation supplies
such witnesses.

For (D.TQ.4), introduce its direct-sum map and
require the stated factorizations of the
natural maps from the threshold transforms.
The quotient maps are surjective, so this
forces the canonical map, rather than an
arbitrary choice.  Impose that it be an
isomorphism.  The Frobenius internal flags
are identified by filtered derived maps;
their one-step Frobenius, fixed normal twists,
and at most $p$ direct sums are fixed
operations.  The chosen shuffle is only a
fixed block pattern.

Finally introduce the forward and inverse
maps between $L_c$ and $F_{a+1}$ with the
fixed exceptional poles $BZ$, and impose
their inverse equations after clearing
$2BZ$.  Require their restrictions to each
$C_i$ to preserve the transformed and target
flags.  Use the coherent six-node tensor
construction of the preceding quadratic
theorem for the pole twists, so that second
twists and compositions are connected to
the original arrows.

Lemmas~\ref{D:polynomial-monad-diagrams} and
\ref{D:filtered-polynomial-diagrams} express
all these maps and equalities in $O(r^2)$
variables and equations of bounded degree
per stage.  In the filtered case there is
one block-triangular Hom space, not a
separate full-rank map for every flag step.
The finite number of threshold transforms
and digit blocks is bounded by $p$, which
is fixed.  All cone, bundle, slope-window
and exactness conditions pass to an open
subset and need no additional quantitative
degree estimate.

This construction is effective over every
local Artinian parameter algebra.  The
relative truncation conditions supply actual
bundles and strict flags.  Locally the
elementary transforms are diagonal with
entries $t$ or $1$; their Frobenius pullbacks
have entries $t^p$ or $1$.  Since $t$ remains
a non-zero-divisor after any parameter base
change, these transformations have exactly
the required images and quotients even
when the parameter ring is nonreduced.
The cone conditions retain any other
negative Tor terms in the ambient
projective-space computations, as in the
preceding quadratic theorem.

\emph{Valid words and class fibres.}
For a length-$f$ type word, recover the unique
weights from $\alpha_{\rm next}=p\alpha-c$.
Discard the invalid words before making any
unit-class assertion.  For a valid word
the common denominator $n$ is prime to $p$.
The exact local ceiling criterion identifies
every witness cycle on $Y^\circ$ with a
unit $p^f$-Frobenius bundle on
$\sqrt[n]{(Y^\circ,C\cap Y^\circ)}$.
Its fixed-frame torsor is finite etale,
and therefore gives a representation of
the original-boundary tame quotient.
Conversely every representation is supplied
by some word and some witness scheme.

Fix a closed witness point and its reference
unit cycle.  The functor of isomorphisms to
the variable cycle is represented using
proper Hom spaces between the underlying
$F_a$, with maps in both directions having
poles at most $BZ$.  Preservation of the
full flags and compatibility with the
one-step arrows are closed equations.
It is an affine functor of finite
presentation, by the same proper Hom
theorem used in the preceding quadratic
proof.

For a small extension of Artinian parameter
rings, a specified unit-isomorphism lifts
uniquely on the root stack, by invariance
of finite etale covers under nilpotent
thickenings.  This can be checked on its
etale Kummer atlas and the corresponding
descent diagram.  Pushforward to the
ordinary parabolic lattices gives a
specified flag-preserving isomorphism
on $Y^\circ$.  The uniform exceptional
pole lemma extends it and its inverse
to $Y$ with the same $B$, independently
of the Artinian length.  Preservation
of flags persists at the points of
$C_i\cap Z$: their quotients are locally
free on $C_i$ over the parameter ring,
so multiplication by a local parameter
of $Z|_{C_i}$ is injective.  A map
which vanishes off those points is
therefore zero.
All cleared inverse and Frobenius
equations extend for the same reason.
The proper unit-Isom scheme thus lifts
the specified Artinian point.  The
infinitesimal smoothness criterion
makes its image open.

Distinct unit classes give disjoint
open subsets of the witness scheme.
They cover its closed points and hence
the whole finite-type scheme.  Thus
there is a finite clopen partition
by representation classes, and their
number is at most the number of
irreducible components.

\emph{The count.}
For each word and multiplicity array,
the witness scheme is an open subset
of a closed affine scheme defined
by $O(fr^2)$ equations of a fixed
degree $D$.  Its number of irreducible
components is at most $D^{O(fr^2)}$.
The surface monad arrays have only
polynomially many choices per stage.
The boundary grade arrays contribute
at most $\exp(Cr)$ per stage by
Lemma~\ref{D:filtered-linear-curve-resolution}.
The ranks are compositions of $r$,
the digit choices number at most
$p^{Cr}$, and the grade permutations
number at most $(r!)^C$, for a fixed
number of boundary curves.
Over $f$ stages their total is
$\exp(O(fr\log(r+1)))$, which is
absorbed by $\exp(O(fr^2))$.
Summing over all valid choices proves
(D.TQ.1).  The final passage to the
original group is by its quotient
map from the group with enlarged
finite omitted set.
\end{proof}
